\documentclass{phd}
\usepackage[utf8]{inputenc}
\usepackage[T1]{fontenc}
\usepackage{geometry}
\usepackage{graphicx}
\usepackage{float}
\usepackage{hyperref}
\usepackage[english,indonesian]{babel}

\title{Optimasi Pemrosesan Template LaTeX Menggunakan Pendekatan Berbasis AI}
\author{Academic Support Agent}
\date{\today}

\begin{document}
\maketitle

\begin{abstract}

\end{abstract}

\chapter{PENDAHULUAN}


\chapter{TINJAUAN PUSTAKA}


% Original content (commented out):
% \documentclass{phd}
%\documentclass[externalize]{phd}
%\documentclass[debug, externalize]{phd}
%\documentclass[debug]{phd}
%\documentclass[draft]{phd}

%\marginnote{Use me to produce a margin footnote.}
%\comment{Use me to make a comment.} <=> \comment[c]{Use me to make a comment.}
%\comment[w]{Use me to make a comment (emphasised as a warning).}
%\comment[e]{Use me to make a comment (emphasised as an error).}

\begin{document}
	%
	\input{./Resources/titlepage}
	%
	\preface
		%
		\input{./Resources/acknowledgements}
		%
		\abstracten
		\abstractfr
		%
		\tableofcontents
		%
	%
	\content
		%
		\input{./Resources/Chapters/Introduction}
		%
		\input{./Resources/Chapters/ChapImportant}
		%
		\input{./Resources/Chapters/ChapVeryImportant}
		%
		\input{./Resources/Chapters/ChapReallyImportant}
		%
		\input{./Resources/Chapters/ChapTheMostImportant}
		%
		\input{./Resources/Chapters/Conclusions}
		%
	%
	\appendix
		%
		\input{./Resources/Appendices/App1}
		%
		\input{./Resources/Appendices/AppVeryImportant}
		%
		\input{./Resources/Appendices/AppTheMostImportant}
		%
		\nocite{*}
		\bibliography{./Resources/Bibliography/Bibliography}
		%
		\listoffigures
		%
	%
	\abstractpage
	%
\end{document}


\end{document}